\begin{enumerate}
	\item Se consolidó el maestro de materiales en un repositorio único, normalizado y trazable. El pipeline de integración construido sobre PostgreSQL con arquitectura medallón unificó los 45,397 registros dispersos en trece archivos de exportación de SAP, resolviendo durante la carga las llaves foráneas contra los catálogos de clases, tipos de material, unidades de medida y UNSPSC. La separación en esquemas de responsabilidad diferenciada permitió que la trazabilidad no fuera un añadido posterior sino una propiedad estructural: toda ejecución de ingesta queda registrada con su archivo de origen, el número de filas procesadas, el estado resultante y el tiempo transcurrido, de modo que cualquier registro del maestro puede auditarse hasta su fuente. El repositorio resultante es el insumo común de los tres módulos analíticos y sustituyó la consulta dispersa sobre hojas de cálculo que caracterizaba la operación previa.

	\item El análisis del catálogo confirmó defectos de calidad de magnitud relevante y aportó un hallazgo no anticipado sobre la taxonomía. La caracterización del maestro reveló que 5,368 registros carecían por completo de categoría asignada en la fuente, y que la distribución de materiales entre clases es fuertemente asimétrica: más de la mitad de ellas cuenta con diez materiales o menos. El análisis de calidad por tipo de material identificó además 2,089 duplicados exactos y 3,649 descripciones que no siguen la convención de separadores del catálogo. Este último dato resultó determinante para el diseño de la solución: sin esa estructura, una descripción es indistinguible de otra similar mediante una búsqueda exacta en SAP, lo que sustentó resolver la detección de duplicados por similitud difusa y no por coincidencia de cadenas. El hallazgo de mayor alcance, sin embargo, no estaba previsto en el planteamiento: el examen de las confusiones del clasificador mostró que una fracción de sus errores no corresponde a fallas del modelo sino a pares de clases que designan el mismo concepto y difieren únicamente en puntuación o separador. La herramienta construida para clasificar materiales resultó también un instrumento de diagnóstico de la taxonomía que los clasifica, capaz de señalar los pares de clases candidatos a fusión.

	\item La configuración más simple de las siete evaluadas resultó también la más exacta, y por un margen consistente. La máquina de vectores de soporte lineal sobre TF-IDF de n-gramas de carácter obtuvo la mejor marca en las cuatro métricas de desempeño (exactitud de 0.8491, F1 macro de 0.7523, F1 ponderado de 0.8380 y exactitud sobre las tres primeras sugerencias de 0.9404) superando a la regresión logística, al bosque aleatorio, a la potenciación de gradiente, al clasificador de subpalabras y al transformador multilingüe preentrenado. La configuración ganadora se entrena en 62.9 segundos sobre procesador convencional, frente a las más de tres horas de la potenciación de gradiente y a los 3,147 segundos del transformador, ambos con resultados inferiores. El modelo seleccionado no solo acierta más sino que se comporta de forma más uniforme a lo largo de una taxonomía en la que más de la mitad de las clases cuenta con diez materiales o menos.

	\item GAMMA integra los tres módulos analíticos en el flujo de creación y mantiene la decisión final en manos del gestor. La plataforma encadena la normalización de la descripción mediante modelo de lenguaje, la búsqueda de duplicados por similitud de trigramas sobre el maestro consolidado y la sugerencia de categoría del clasificador, y los presenta durante la creación del material y no en un reporte posterior. Como cada sugerencia lleva una probabilidad asociada, la plataforma resuelve por sí sola los casos en que esa probabilidad supera el umbral configurado —el 42\,\% del volumen, con una exactitud del 98.6\,\%— y deriva el resto a revisión. Ninguna sugerencia llega a SAP sin la confirmación del gestor: el sistema asiste la decisión, no la sustituye. El comportamiento en operación respaldó ese diseño: sobre las 210 solicitudes asistidas durante el período de validación, el gestor conservó la categoría sugerida en el 89.0\,\% de los casos. La encuesta de cierre apuntó en la misma dirección; los gestores destacaron la facilidad de uso y la jefatura la alineación con los objetivos de calidad del área, mientras que la precisión de la categoría sugerida fue el aspecto peor valorado.

	\item La herramienta redujo a menos de la mitad el tiempo de gestión, con un retorno de la inversión inmediato. El tiempo efectivo del gestor pasó de 3.34 a 1.39 horas por solicitud, una reducción del 58\,\% que las pruebas estadísticas confirman como significativa y no atribuible al azar ni al menor tamaño del período de validación. La reducción tampoco se explica por una carga de trabajo menor: durante ese mismo período el equipo atendió más solicitudes por día hábil que antes de la implementación. Bajo el escenario conservador, el ahorro proyectado asciende a Q400,491 anuales frente a un costo total del proyecto de Q1,641.65, de modo que la inversión se recupera en cerca de un día y medio de operación.
\end{enumerate}

En cuanto al objetivo general, GAMMA se desarrolló, se desplegó y operó sobre solicitudes reales del proceso de creación individual de materiales, articulando efectivamente el procesamiento de lenguaje natural, la búsqueda por similitud difusa y el aprendizaje automático supervisado sobre el histórico del catálogo. La reducción del tiempo de gestión quedó demostrada con significancia estadística sobre la operación real, y el retorno de la inversión resulta inmediato bajo cualquiera de los escenarios considerados.

El aporte a la calidad del maestro quedó igualmente respaldado, aunque por una vía distinta. El acuerdo del 89.0\,\% entre la sugerencia del modelo y la decisión final del gestor constituye una validación del clasificador independiente de su evaluación de laboratorio: mide su desempeño contra las decisiones efectivas de quienes conocen el catálogo, y no contra una partición retenida de los mismos datos con que fue entrenado. Que ambas mediciones converjan —84.91\,\% sobre datos retenidos y 89.0\,\% sobre la operación— indica que el desempeño estimado en el laboratorio se sostuvo al enfrentar materiales reales.

Resta, no obstante, una distinción que conviene declarar. Que el modelo acierte en operación y que los tres controles intervengan antes de cada alta establece que la herramienta previene la incorporación de nuevos defectos; no establece todavía en qué medida el catálogo mejora como consecuencia. Esa segunda afirmación exige contrastar la evolución de los indicadores de duplicidad y de categorización sobre un horizonte de operación más extenso que las cinco semanas del piloto, y constituye la continuación natural del trabajo.
